% !TEX root = ./00-main.tex

\chap{Abstract} 

How can cheap model predictions be combined with a much smaller set of gold-standard labels to support statistically reliable estimation and uncertainty quantification? Prediction-powered inference offers one answer: it uses predictions available at scale together with a smaller labeled subset to estimate the quantity of interest and quantify its uncertainty. This thesis develops that idea in two methodological directions and one applied case study. The first foundation asks how prediction-powered correction behaves when the cheap signal is itself noisy, as in large-language-model (LLM) judge labels. The second asks how surrogate quality changes the number of gold-standard labels needed for a target power calculation.

The applied case study focuses on physical protein--protein interactions, meaning whether two proteins are recorded as interacting in curated benchmark data. I develop a leakage-aware workflow that builds edge-level surrogate scores from protein sequence and UniProt text representations and then carries a fixed ensemble predictor into downstream estimation and planning. The predictive goal is to score candidate interactions. The inferential goal is different: estimate prevalence-like quantities and their uncertainty when only a limited subset of edges is experimentally validated.

Across the main held-out protein benchmarks, prediction-powered correction remains close to the corresponding full-data reference values even when only a small labeled subset is used. In the planning layer, stronger surrogates reduce the labeled sample size needed for mean- or proportion-type targets, leading to moderate gains. Overall, the protein application shows that a carefully constructed surrogate can support corrected estimation and practical planning, even when the resulting gains are limited by the strength of the available prediction model.

\par\medskip


%% list of keywords seperated by comma
\keywords{Biostatistics, prediction-powered inference, surrogate-assisted inference, LLM-as-a-judge, protein--protein interaction, power analysis}


%%%%  committee members (add it right after the abstract w/o page break)
\begin{singlespace}

    %% if you have co-advisor, add here w/ \vspace{0.1in} as shown below
    %% alternatively you can use minipage environment to put side-by-side
    \section*{Primary reader and thesis advisor}
    
    Dr. Yiqun Chen \\
    Assistant Professor\\
    Department of Biostatistics\\
    Johns Hopkins University Bloomberg School of Public Health, Baltimore MD 


    \section*{Secondary reader}
    
    Dr. Ingo Ruczinski\\
    Professor\\
    Department of Biostatistics \\
    Johns Hopkins University Bloomberg School of Public Health, Baltimore, MD 
    
    %% you can add more readers if you have them on your committee 
    %% use \vspace{0.1in} in between members for clarity
    %% you can also place committee members side-by-side using `minipage`

\end{singlespace}
